%%==========================================================================%%
%% abstract.tex -- working abstract for the WBPET pediatric whole-body       %%
%% PET/MRI dataset descriptor (Nature Scientific Data).                      %%
%%                                                                           %%
%% NOTE: Springer Nature requires single-file submission. Before submitting, %%
%% inline the contents of this file into main.tex (replacing %%==========================================================================%%
%% abstract.tex -- working abstract for the WBPET pediatric whole-body       %%
%% PET/MRI dataset descriptor (Nature Scientific Data).                      %%
%%                                                                           %%
%% NOTE: Springer Nature requires single-file submission. Before submitting, %%
%% inline the contents of this file into main.tex (replacing %%==========================================================================%%
%% abstract.tex -- working abstract for the WBPET pediatric whole-body       %%
%% PET/MRI dataset descriptor (Nature Scientific Data).                      %%
%%                                                                           %%
%% NOTE: Springer Nature requires single-file submission. Before submitting, %%
%% inline the contents of this file into main.tex (replacing %%==========================================================================%%
%% abstract.tex -- working abstract for the WBPET pediatric whole-body       %%
%% PET/MRI dataset descriptor (Nature Scientific Data).                      %%
%%                                                                           %%
%% NOTE: Springer Nature requires single-file submission. Before submitting, %%
%% inline the contents of this file into main.tex (replacing \input{abstract}%%
%% with the \abstract{...} block below).                                     %%
%%==========================================================================%%

\abstract{Recently, image-based AI tools have shown great promise across clinical applications,
and a key upstream task is the accurate identification of anatomy, against which metabolically active lesions can be localized.
Although several datasets support whole-body anatomical analysis, to date they have been limited to adults,
which has hampered pediatric adoption, since models calibrated on adults perform unpredictably on pediatric inputs.
Adult whole-body CT datasets are also of limited use,
because MRI is preferred in children to reduce radiation exposure, and anatomy alone does not capture disease activity.
What is needed, therefore, is a whole-body pediatric PET/MRI dataset that addresses these entangled gaps in coverage, modality, and functional labeling.
We therefore present a pediatric whole-body PET/MRI dataset of [TBD] examinations, spanning ages [TBD],
with active lymphoma and disease-free follow-up examinations and their anatomical and tumor-lesion segmentations.
We describe the dataset's acquisition, organization, and validation,
and discuss its use for developing and benchmarking pediatric whole-body anatomical and oncological analysis tools.}
%%
%% with the \abstract{...} block below).                                     %%
%%==========================================================================%%

\abstract{Recently, image-based AI tools have shown great promise across clinical applications,
and a key upstream task is the accurate identification of anatomy, against which metabolically active lesions can be localized.
Although several datasets support whole-body anatomical analysis, to date they have been limited to adults,
which has hampered pediatric adoption, since models calibrated on adults perform unpredictably on pediatric inputs.
Adult whole-body CT datasets are also of limited use,
because MRI is preferred in children to reduce radiation exposure, and anatomy alone does not capture disease activity.
What is needed, therefore, is a whole-body pediatric PET/MRI dataset that addresses these entangled gaps in coverage, modality, and functional labeling.
We therefore present a pediatric whole-body PET/MRI dataset of [TBD] examinations, spanning ages [TBD],
with active lymphoma and disease-free follow-up examinations and their anatomical and tumor-lesion segmentations.
We describe the dataset's acquisition, organization, and validation,
and discuss its use for developing and benchmarking pediatric whole-body anatomical and oncological analysis tools.}
%%
%% with the \abstract{...} block below).                                     %%
%%==========================================================================%%

\abstract{Recently, image-based AI tools have shown great promise across clinical applications,
and a key upstream task is the accurate identification of anatomy, against which metabolically active lesions can be localized.
Although several datasets support whole-body anatomical analysis, to date they have been limited to adults,
which has hampered pediatric adoption, since models calibrated on adults perform unpredictably on pediatric inputs.
Adult whole-body CT datasets are also of limited use,
because MRI is preferred in children to reduce radiation exposure, and anatomy alone does not capture disease activity.
What is needed, therefore, is a whole-body pediatric PET/MRI dataset that addresses these entangled gaps in coverage, modality, and functional labeling.
We therefore present a pediatric whole-body PET/MRI dataset of [TBD] examinations, spanning ages [TBD],
with active lymphoma and disease-free follow-up examinations and their anatomical and tumor-lesion segmentations.
We describe the dataset's acquisition, organization, and validation,
and discuss its use for developing and benchmarking pediatric whole-body anatomical and oncological analysis tools.}
%%
%% with the \abstract{...} block below).                                     %%
%%==========================================================================%%

\abstract{Recently, image-based AI tools have shown great promise across clinical applications,
and a key upstream task is the accurate identification of anatomy, against which metabolically active lesions can be localized.
Although several datasets support whole-body anatomical analysis, to date they have been limited to adults,
which has hampered pediatric adoption, since models calibrated on adults perform unpredictably on pediatric inputs.
Adult whole-body CT datasets are also of limited use,
because MRI is preferred in children to reduce radiation exposure, and anatomy alone does not capture disease activity.
What is needed, therefore, is a whole-body pediatric PET/MRI dataset that addresses these entangled gaps in coverage, modality, and functional labeling.
We therefore present a pediatric whole-body PET/MRI dataset of [TBD] examinations, spanning ages [TBD],
with active lymphoma and disease-free follow-up examinations and their anatomical and tumor-lesion segmentations.
We describe the dataset's acquisition, organization, and validation,
and discuss its use for developing and benchmarking pediatric whole-body anatomical and oncological analysis tools.}
